problem:  
Let \((M,g)\) be a complete, noncompact Riemannian manifold with nonnegative Ricci curvature. Suppose there exists a nonconstant harmonic function \( f : M \to \mathbb{R} \) such that:  
1. every regular level set \( f^{-1}(t) \) is a compact minimal hypersurface, and  
2. the gradient \( \nabla f \) satisfies \( |\nabla f|(x) \to c > 0 \) uniformly as \( x \to \infty \) (i.e., outside every compact set there exists \(R>0\) s.t. \(||\nabla f|(x) - c| < \varepsilon\) whenever \(d(x,o) > R\)).  

Does \( (M,g) \) admit an **asymptotic splitting** in the sense that there exists a sequence of basepoints \(x_i \to \infty\) for which the pointed manifolds \((M,g,x_i)\) converge in the pointed Gromov–Hausdorff topology to a product space \( (N \times \mathbb{R}, g_N + ds^2) \)?  

Why is it a "good" problem:  
This problem asks whether analytic data from a harmonic function—specifically, compact minimal level sets and asymptotically constant gradient—can enforce a geometric splitting at infinity under nonnegative Ricci curvature. It generalizes the spirit of the Cheeger–Gromoll splitting theorem by replacing the explicit existence of a geodesic line with the asymptotic analytic behavior of a potential function. The formulation is simple, yet it sits at the cutting edge of geometric analysis: it connects harmonic function theory, minimal hypersurface geometry, and the large-scale structure of manifolds with Ricci \(\ge 0\). Establishing such an asymptotic product structure (or finding counterexamples) would likely deepen the analytic understanding of lines, potential functions, and curvature rigidity, thus bridging elliptic PDE methods and global Riemannian geometry.